\begin{textsample}{NEG Dim 2 – 2010\_c4 – Score 5.00 – 2010\_c4/enade\_2010\_tec\_gest\_hosp\_1a50307899.txt}  \label{ex:f2_neg_287}
CAPÍTULO 4 PERCEPÇÃO DA PROVA As análises feitas neste capítulo tratam das percepções dos estudantes, separados por grupos de Ingressantes e Concluintes, sobre a prova aplicada no ENADE/2010 na Área de Tecnologia em Gestão Hospitalar.
Estas percepções foram mensuradas por meio de nove questões que avaliaram desde o grau de dificuldade da prova até o tempo gasto para concluí -la.
As percepções sobre a prova foram relacionadas com o desempenho dos alunos e com a Grande Região de funcionamento do curso.
O desempenho dos alunos foi classificado em dois níveis : quarto inferior e quarto superior.
Para tanto, esse desempenho foi colocado em ordem ascendente.
O percentil 25, P25, também conhecido como primeiro quartil, é a nota de desempenho que deixa um quarto ( 25,0 \% ) dos valores observados abaixo e três quartos acima.
A Figura 1 apresenta uma ilustração destes conceitos.
O quarto inferior de desempenho é composto pelas notas abaixo do primeiro quartil.
Já o percentil 75, P75, também conhecido como terceiro quartil, é o valor para o qual há três quartos ( 75,0 \% ) dos dados abaixo e um quarto acima dele.
O quarto superior de desempenho é composto pelas notas iguais ou acima do terceiro quartil.
O percentil 50, P50, também conhecido como mediana, é o valor que divide as notas em dois conjuntos de igual tamanho.
Vale ressaltar que percentis, quartis e medianas são pontos, e não obrigatoriamente pertencentes ao conjunto original de dados.
Ao passo que os quartos são subconjuntos dos dados originais.
A seguir, serão apresentados gráficos com resultados selecionados, relativos às nove questões avaliadas por grupos de estudantes.
As Tabelas no Anexo II apresentam os valores absolutos e a distribuição percentual das alternativas válidas das nove questões, segundo o mesmo recorte de desempenho dos alunos e Grande Região de funcionamento do curso. 4.1 GRAU DE DIFICULDADE DA PROVA 4.1.1 Componente de Formação Geral Na questão “ Qual o grau de dificuldade desta prova na parte de Formação Geral? ”, 25,8 \% do grupo de Ingressantes e 20,9 \% do grupo de Concluintes optaram pelas alternativas difícil ou muito difícil.
Entretanto, para mais da metade dos estudantes, tanto Ingressantes ( 65,4 \% ) quanto Concluintes ( 72,4 \% ), o Componente de Formação Geral da prova foi considerado com grau de dificuldade médio ( Gráfico 4.1 e Tabelas II.1 e II.10 no Anexo II ).
O percentual de Ingressantes que consideraram a prova difícil ou muito difícil foi mais elevado no quarto inferior ( 34,9 \% ) do que no superior ( 22,3 \% ), o que também foi verificado no grupo dos Concluintes, com percentual de participantes maior no quarto inferior ( 28,0 \% ), do que no quarto superior ( 13,2 \% ).
Semelhantemente ao ocorrido para os Ingressantes como um todo, para os quartos de menor e maior desempenho deste grupo, a alternativa modal para esta pergunta foi médio, com respectivamente 56,5 \% e 67,3 \% dos respondentes.
Já entre os Concluintes, estas percentagens foram, respectivamente, 69,3 \% e 76,8 \%.
Considerando -se as alternativas escolhidas pelos estudantes de cada Grande Região, a região Centro-Oeste foi aquela em que houve maior proporção de Ingressantes que avaliaram o Componente de Formação Geral da prova como difícil ou muito difícil : 34,5 \% tiveram essa percepção.
Os Concluintes, por sua vez, emitiram essa opinião em maior medida na região Sudeste ( 27,3 \% ).
Nas Grandes Regiões, a proporção maior de Ingressantes que consideraram o Componente de Formação Geral da prova como sendo de grau de dificuldade médio esteve localizada entre 60,1 \% na região Centro-Oeste e 77,3 \% na região Norte.
Optando pela alternativa de grau médio por Grande Região, no grupo dos Concluintes os percentuais estiveram entre 69,1 \% na região Sudeste e 74,8 \%, na Nordeste. 4.1.2 Componente de Conhecimento Específico Ao responderem à questão “ Qual o grau de dificuldade desta prova na parte de Componente Específico? ”, 36,6 \% do grupo de Ingressantes e 28,0 \% do grupo de Concluintes classificaram -na como difícil ou muito difícil.
Este Componente de Conhecimento Específico da prova foi considerado com grau de dificuldade médio por 56,6 \% dos Ingressantes e por 64,5 \% dos Concluintes ( Gráfico 4.2 e Tabelas II.2 e II.11 no Anexo II ).
Considerando -se a percepção sobre o grau de dificuldade do Componente de Conhecimento Específico da prova de acordo com o desempenho dos estudantes, observa -se que 38,5 \% dos Ingressantes obtiveram desempenho situado no quarto superior classificaram -na como difícil ou muito difícil.
Dos Ingressantes que obtiveram desempenho situado no quarto inferior, 35,8 \% optaram pela mesma alternativa.
Para os Concluintes, o percentual de estudantes que consideraram o Componente de Conhecimento Específico da prova difícil ou muito difícil foi, no quarto superior, 31,6 \%, e 27,8 \% no quarto inferior.
Para os quartos de menor e maior desempenho, a alternativa modal para esta pergunta, tanto para Ingressantes quanto para Concluintes, foi médio : 55,1 \% dos Ingressantes do quarto superior e 55,5 \% do quarto inferior deram essa resposta.
Quanto aos Concluintes, 62,6 \% do quarto superior e 63,5 \% do inferior assim se pronunciaram.
Entre as Grandes Regiões brasileiras, a proporção de Ingressantes que avaliaram o Componente de Conhecimento Específico da prova como difícil ou muito difícil variou de 31,9 \% ( Norte ) a 41,5 \% ( Centro-Oeste ).
O grupo de Concluintes avaliou esse componente como difícil ou muito difícil com menor frequência, com o mínimo alcançado na região Norte, 23,2 \%, e o máximo na região Nordeste, com 29,8 \% dos respondentes.
O percentual de Ingressantes que classificaram o grau de dificuldade como médio, no Componente de Conhecimento Específico da prova, variou de 52,1 \% ( Centro-Oeste ) a 61,0 \% ( Norte ).
Entre os Concluintes, a percepção do mesmo grau de dificuldade, variou de 60,3 \% ( Sul ) a 68,2 \% ( Norte ). 4.2 EXTENSÃO DA PROVA EM RELAÇÃO AO TEMPO TOTAL Indagados quanto à extensão da prova em relação ao tempo total oferecido para a sua resolução, os estudantes apontaram com maior incidência a alternativa que considerava a extensão adequada, para ambos os grupos e todas as agregações consideradas ( Gráfico 4.3 e Tabelas II.3 e II.12 no Anexo II ).
No grupo de Ingressantes, a maioria, 61,5 \%, respondeu ser a extensão da prova adequada, ao passo que, no grupo de Concluintes, 60,7 \% tiveram a mesma opinião.
Para 34,2 \% dos Ingressantes o exame foi longo ou muito longo.
Essa mesma consideração foi feita por 34,8 \% dos Concluintes.
No grupo de Ingressantes, a maioria, tanto para o quarto de desempenho inferior ( 57,9 \% ) quanto para o de melhor desempenho ( 64,6 \% ) respondeu ser a extensão da prova adequada.
Entre os Concluintes, esses valores foram situados em cada quarto em, respectivamente, 57,5 \% e 61,9 \%.
Considerando -se ainda os quartos extremos de desempenho, observa -se que a proporção de estudantes que avaliaram a prova longa ou muito longa, em relação ao tempo total destinado à sua resolução, acompanhou a percepção de Ingressantes e Concluintes.
Para o grupo de Ingressantes, entre os estudantes no quarto de desempenho inferior, 36,7 \% declararam ter esta percepção, ao passo que, entre aqueles no quarto de melhor desempenho, esta proporção foi de 32,2 \%.
O mesmo pode ser constatado entre os Concluintes : 63,5 \% no quarto de menor desempenho e 33,9 \% no quarto de maior desempenho declararam ter tido esta percepção.
Na análise por Grandes Regiões, o percentual de Ingressantes que responderam ser a extensão da prova adequada variou de 60,0 \% no Nordeste a 72,7 \% no Sul.
Já entre os Concluintes, estes valores estavam compreendidos entre 56,3 \% ( Sudeste ) e 67,7 \% ( Norte ).
Aqueles Ingressantes que avaliaram a prova como longa ou muito longa em relação ao tempo total destinado à sua resolução variou de 23,4 \% na região Sul a 35,8 \% na região Centro-Oeste.
Os percentuais, quanto ao mesmo aspecto, entre os Concluintes, foram de 26,7 \% na região Norte, indo até 38,2 \% na região Sudeste. 4.3 COMPREENSÃO DOS ENUNCIADOS DAS QUESTÕES 4.3.1 Componente de Formação Geral Com relação aos enunciados das questões do Componente de Formação Geral, as opiniões foram positivas, com percentuais que foram dos 72,9 \% para Concluintes aos 72,0 \% para Ingressantes que consideraram claros e objetivos os enunciados de todas ou da maioria das questões ( Gráfico 4.4 e Tabelas II.4 e II.13 no Anexo II ).
Ao se comparar as percepções dos estudantes segundo o desempenho, observou -se que, entre os Ingressantes do quarto superior, a clareza e objetividade de todos ou da maioria dos enunciados das questões foram percebidas por 79,6 \% dos estudantes.
A proporção apurada para os Concluintes desse mesmo quarto foi 80,2 \%.
No quarto de desempenho inferior, tal avaliação foi emitida por 62,5 \% dos Ingressantes e por 56,2 \% dos Concluintes.
A partir dessas estatísticas, é possível afirmar que todos ou a maioria dos enunciados das questões do Componente de Formação Geral estavam claros e objetivos para a maioria dos respondentes, segundo os quartos de desempenho considerados, embora os que se situaram no quarto superior tenham declarado esta percepção com maior incidência.
Na análise regional, observa -se que as opiniões de que todos ou a maioria dos enunciados das questões do Componente de Formação Geral estavam claros e objetivos variaram, entre os Ingressantes, de 68,1 \% na região Sudeste a 89,6 \% na região Sul.
Entre os Concluintes, os valores correspondentes variaram de 64,5 \% no Centro-Oeste a 78,3 \%, no Sul. 4.3.2 Componente de Conhecimento Específico Com relação aos enunciados das questões do Componente de Conhecimento Específico da prova, para os Concluintes da Área de Tecnologia em Gestão Hospitalar, a clareza e a objetividade foram percebidas em todos ou na maioria destes por 71,7 \%.
Para o grupo dos Ingressantes, foram percebidas por 67,2 \%.
A proporção de estudantes que considerou os enunciados das questões claros e objetivos foi mais elevada nos quartos superiores : 73,5 \% dos Ingressantes e 78,5 \% dos Concluintes, se comparados aos 55,0 \% e aos 57,5 \% dos mesmos grupos, respectivamente, que estiveram no quarto inferior de desempenho ( Gráfico 4.5 e Tabelas II.5 e II.14 no Anexo II ).
Nas cinco Grandes Regiões brasileiras, os percentuais de Concluintes que consideraram claros e objetivos todos ou a maioria dos enunciados das questões do Componente de Conhecimento Específico da prova foram semelhantes aos de Ingressantes.
Os Ingressantes ( 86,9 \% ) e os Concluintes ( 81,9 \% ) da região Sul apresentaram as proporções mais elevadas desta percepção.
Os Ingressantes da região Sudeste ( 61,8 \% ) e os Concluintes da região Centro-Oeste ( 66,2 \% ) apresentaram as proporções mais baixas dessa percepção. 4.4 SUFICIÊNCIA DAS INFORMAÇÕES/INSTRUÇÕES FORNECIDAS Ao avaliarem as informações/instruções fornecidas para a resolução das questões, 71,6 \% do grupo dos Ingressantes e 79,8 \% do grupo dos Concluintes respondentes de todo o Brasil afirmaram que estas eram até excessivas ou suficientes em todas ou na maioria das questões ( Gráfico 4.6 e Tabelas II.6 e II.15 no Anexo II ).
Levando -se em conta o desempenho dos participantes, notam -se diferenças entre as opiniões de estudantes, Ingressantes e Concluintes, dos quartos inferior e superior de desempenho, como mostra o Gráfico 4.6.
O percentual de participantes que avaliaram as informações/instruções como excessivas ou suficientes em todas ou na maioria das questões foi mais elevado nos quartos superiores, correspondendo a 82,5 \% dos Ingressantes e a 86,8 \% dos Concluintes, percentuais superiores à média nacional.
Já nos quartos inferiores, a suficiência das informações/instruções declarada como excessiva em todas ou na maioria das questões foi percebida por 59,3 \% dos Ingressantes e por 71,3 \% dos Concluintes.
Quanto à distribuição de respondentes nas Grandes Regiões, observa -se que as proporções de estudantes que perceberam como excessivas ou suficientes em todas ou na maioria das questões as informações/instruções fornecidas para a resolução das mesmas superaram as médias nacionais, em maior medida, entre os Ingressantes da região Sul ( 84,4 \% ).
Entre os Concluintes, a superação do total nacional foi mais expressiva na região Norte ( 83,1 \% ). 4.5 DIFICULDADE ENCONTRADA AO RESPONDER À PROVA Perguntados sobre as dificuldades com as quais se depararam ao responder à prova, 30,1 \% dos Ingressantes e 14,6 \% dos Concluintes, apontaram o desconhecimento do conteúdo.
Concluintes, por sua vez, apontaram, com maior frequência, a forma diferente de abordagem do conteúdo como dificuldade, mencionada por 51,1 \%.
Esta alternativa foi mencionada por uma parcela menor dos Ingressantes, 42,2 \%.
A falta de motivação para fazer a prova foi a dificuldade para 16,7 \% dos Concluintes e para 13,0 \% dos Ingressantes.
Considerando -se os respondentes em todo o Brasil, 12,8 \% dos Ingressantes e 15,1 \% dos Concluintes afirmaram que não tiveram qualquer tipo de dificuldade para responder à prova ( Tabelas II.7 e II.16 no Anexo II ).
Com relação aos quartos de desempenho, a forma diferente de abordagem do conteúdo foi a dificuldade modal para Ingressantes e Concluintes.
Foi apontada por 43,4 \% daqueles no quarto superior, entre Ingressantes, e por 39,9 \% dos do quarto inferior.
Entre os Concluintes, os percentuais de estudantes com a mesma percepção foi 49,6 \% no quarto superior e 48,3 \% no quarto inferior.
Na análise por Grandes Regiões, o percentual de Ingressantes que apontaram o desconhecimento do conteúdo como dificuldade ao responder à prova variou de 17,3 \%, na região Norte a 42,7 \%, na Sudeste.
Entre os Concluintes, tal percentual foi mais baixo, variando de 7,0 \% na região Norte a 20,9 \%, na Nordeste.
A forma diferente de abordagem do conteúdo foi citada, entre os Concluintes, por percentuais que variaram de 43,8 \% ( região Nordeste ) a 60,3 \% ( Sul ).
Os percentuais relacionados a essa opção entre os Ingressantes variaram de 28,9 \% ( região Sul ) a 47,3 \% ( região Centro-Oeste ).
O percentual de Ingressantes que citou a falta de motivação como dificuldade superou a média nacional, em maior medida, na região Norte ( 18,0 \% ).
Entre os Concluintes, esse percentual variou de 8,4 \% ( região Sul ) a 23,9 \% ( região Norte ).
Houve percentuais acima do total nacional dos Ingressantes nas regiões Nordeste ( 13,3 \% ), Norte ( 16,5 \% ) e Sul ( 22,4 \% ), com declarações de que não tiveram qualquer dificuldade para responder à prova.
Entre os Concluintes, esta situação aconteceu, em maior proporção na região Norte ( 19,7 \% ).
O Gráfico 4.7 apresenta os percentuais de estudantes, Ingressantes e Concluintes, que apontaram o desconhecimento do conteúdo como dificuldade percebida ao responder à prova. 4.6 CONTEÚDOS DAS QUESTÕES OBJETIVAS DA PROVA Ao analisarem os conteúdos das questões objetivas da prova, a maior parte dos Ingressantes, 50,9 \% do total nacional, afirmou que não estudou ainda a maioria desses conteúdos, enquanto entre os Concluintes, 10,5 \% expressaram a mesma opinião.
A maioria dos Concluintes ( 61,4 \% ) afirmou ter estudado e aprendido muitos ou todos os conteúdos.
Para os Ingressantes, o percentual correspondente a essas alternativas foi 26,1 \% ( Gráfico 4.8 e Tabelas II.8 e II.17 no Anexo II ).
Considerando -se separadamente as opiniões de estudantes dos quartos de desempenho inferior e superior, observa -se que, para os Ingressantes, o percentual de participantes que afirmaram que não estudaram ainda a maioria desses conteúdos, foi mais elevado no quarto inferior ( 49,1 \% ) do que no superior ( 46,4 \% ).
Para os Concluintes do quarto inferior, 15,8 \% ofereceram essa resposta, sendo 5,9 \% os do quarto superior.
Tendo em conta os quartos superiores para os dois grupos, 80,6 \% dos Concluintes e 34,4 \% dos Ingressantes afirmaram ter estudado e aprendido muitos ou todos os conteúdos.
Na análise por Grande Região, a proporção de Ingressantes que afirmou que não estudou ainda a maioria desses conteúdos superou o total nacional nas regiões Sudeste ( 53,6 \% ) e Nordeste ( 56,9 \% ).
Já entre os Concluintes, as regiões Nordeste ( 13,1 \% ) e Sudeste ( 15,5 \% ) apresentaram percentuais maiores do que a média nacional para esta alternativa.
Em todas as regiões, Concluintes afirmaram ter estudado e aprendido muitos ou todos os conteúdos, variando entre 49,9 \% na região Sudeste e 74,6 \% na Norte. 4.7 TEMPO GASTO PARA CONCLUIR A PROVA Ao responderem sobre o tempo de conclusão da prova, mais da metade dos estudantes afirmou ter gasto entre duas e quatro horas.
Considerando -se separadamente as percepções do grupo de Ingressantes e do grupo de Concluintes, os percentuais, em termos nacionais, foram de 64,9 \% e 64,1 \%, respectivamente ( Gráfico 4.9 e Tabelas II.9 e II.18 no Anexo II ).
Uma vez sendo analisadas as alternativas oferecidas pelos estudantes que se situaram nos quartos superior e inferior, observa -se que maior proporção de Ingressantes no quarto superior declarou ter gasto entre duas e quatro horas para concluir a prova quando comparados com os do quarto inferior, respectivamente, 74,6 e 52,2 \%.
A mesma situação foi verificada no grupo dos Concluintes, no qual 80,3 \% daqueles no quarto superior e 47,8 \% daqueles no quarto inferior declaram ter gasto entre duas e quatro horas para concluir a prova.
Nas cinco Grandes Regiões brasileiras, os Ingressantes que utilizaram entre duas a quatro horas para finalizar a prova superaram o percentual nacional nas regiões Sul ( 67,5 \% ), Centro-Oeste ( 67,6 \% ) e Norte ( 71,5 \% ).
Entre os Concluintes, o mesmo foi verificado entre os estudantes das regiões Norte ( 65,3 \% ), Sudeste ( 66,7 \% ), Sul ( 71,2 \% ) e Centro-Oeste ( 71,0 \% ).
A região Nordeste foi a que apresentou o menor percentual de Concluintes que dispensaram entre duas e quatro horas para concluir a prova : 55,2 \%, como mostra o Gráfico 4.9.

% matched lemmas: 
\end{textsample}
